\documentclass{article}
\usepackage[T1]{fontenc}
\usepackage{lmodern}
\usepackage{graphicx}
\usepackage{amsmath,amssymb}
\usepackage[a4paper,margin=2cm]{geometry}
\usepackage[absolute,overlay]{textpos}
\setlength{\TPHorizModule}{1mm}
\setlength{\TPVertModule}{1mm}
\usepackage[indonesian]{babel}
\usepackage{hyperref}
\setlength{\emergencystretch}{2em}
% unit: Halaman 12

\begin{document}

% === Halaman 12 (python/native) ===

12

Algoritma Pertukaran Kunci Diffie-Hellman

1.

Alice membangkitan bilangan bulat acak a  dan mengirim hasil perhitungan 

berikut kepada Bob:



A = ga mod p

2.

Bob membangkitkan bilangan bulat acak b dan mengirim hasil perhitungan 

berikut kepada Alice:



B = gb mod p

3.

Alice menghitung



K = Ba mod p  


(K  = gab mod p)

4.

Bob menghitung



K = Ab mod p  


(K  = gab mod p)

• Jika perhitungan dilakukan dengan benar, maka  K = gab mod p. Alice dan Bob 

sekarang memiliki kunci rahasia yang sama, yaitu K.

(a = kunci privat Alice, A = kunci publik Alice) 

(b = kunci privat Bob, B = kunci publik Bob) 

Rinaldi M/II4020 Kriptografi

\end{document}
